
@@M0@@
\uline{où} @@M1@@.

On en déduit par récurrence immédiate sur @@M2@@ :

\textbf{Corollaire 1.1.4.} \uline{Pour tout schéma noethérien} @@M3@@ \uline{et pour tout entier} @@M4@@, @@M5@@ \uline{est le sous-groupe de} @@M6@@ \uline{engendré par les classes} @@M7@@\uline{, où} @@M8@@ \uline{parcourt l'ensemble des sous-schémas intègres fermés de dimension} @@M9@@ \uline{de} @@M10@@.

\emph{Démonstration de} 1.1.2 : Soit @@M11@@. Comme @@M12@@ est un espace noethérien de dimension @@M13@@, les éléments de @@M14@@ sont maximaux dans @@M15@@, donc en nombre fini, et @@M16@@ est finie pour tout @@M17@@. Cela prouve la première assertion.

Soit @@M18@@ la catégorie des @@M19@@-modules cohérents et soit @@M20@@ le sous-ensemble de @@M21@@ formé des @@M22@@-modules cohérents concentrés sur @@M23@@ qui satisfont à la condition (1.1.3). Évidemment @@M24@@ contient les faisceaux @@M25@@ avec @@M26@@. Il suffit donc de montrer que @@M27@@ est un sous-ensemble «~exact~» de @@M28@@, car alors le lemme de dévissage (EGA III 3.1.2) donne @@M29@@.

Soit donc
@@M30@@
une suite exacte de @@M31@@-modules ayant deux des membres dans @@M32@@. Alors les faisceaux @@M33@@, @@M34@@ et @@M35@@ sont tous concentrés sur @@M36@@, donc les longueurs @@M37@@, @@M38@@ et @@M39@@ des @@M40@@-modules @@M41@@, @@M42@@ et @@M43@@ sont finies
